\section{Discussion}
\label{sec:discussion}

The central empirical finding is that a patient-blocked alarm budget is portable as a \emph{constraint} but not as a \emph{threshold}. Within site A, the calibrated policy moves smoothly from 0.19 sensitivity at 0.25 alarms per patient-day to 0.77 at 2 per day, and the realised rates remain below the requested budgets. When the same A-calibrated threshold is applied to site B, the realised rate falls to roughly half the A rate and sensitivity at 1 per day falls from 0.585 to 0.315. The budget itself is still respected in the sense that B's realised mean is below the A certificate, but the clinical yield is much lower. This distinction matters for engineering practice: a site that imports only a threshold has not imported the operating point that the source site validated.

The comparison with the timestep shortcut clarifies what patient blocking buys. At the tightest budgets the timestep procedure selects a lower threshold and increases both sensitivity and the fraction of patients exceeding the budget. That is a real trade, not a statistical technicality: the method that treats correlated hours as independent observations is more willing to spend a fixed nominal risk budget, and the additional alarms are not free at the patient level. At looser budgets the finite grid collapses the two procedures to the same threshold, so the experiment should not be read as a claim that timestep calibration always fails. The defensible conclusion is narrower and stronger: when the estimand is a patient-average rate, the effective sample size is the number of independent patients, and methods that substitute timesteps can become less conservative without a corresponding patient-level guarantee.

The backbone comparison is more nuanced than a single ranking. LightGBM has the highest site-B AUROC ($0.713\pm0.003$) and the highest external sensitivity at 1 per day ($0.540\pm0.071$), while the GRU has the highest site-A sensitivity ($0.684\pm0.006$) and LightGBM is second ($0.661\pm0.042$). The seed-0 paired LightGBM-minus-logistic sensitivity difference at B is $0.081$ with interval $[0.027,0.133]$ and Holm-adjusted $p=0.066$, whereas the GRU-minus-logistic difference is $-0.056$ with an interval that includes zero after Holm adjustment. A stronger tabular scorer therefore narrows the external gap, but it does not remove the transport problem: LightGBM's B rate is $0.410\pm0.156$, still below the requested 1 per day, and the target label-budget analysis shows that local labels remain necessary to move the operating point. This is consistent with the broader observation that clinical prediction performance is often more sensitive to cohort and workflow differences than to modest architecture changes \citep{valan2025evaluating,ostermayer2024external,wu2026outofdistribution}. It also argues against presenting a backbone-only AUROC table as evidence of deployability.

For engineering deployment, the results support four concrete practices. First, the minimum observed stay used in the support bound should be declared before looking at calibration or test outcomes, because it changes both eligibility and $L_{\max}$; the 24-hour restriction here removes about 22\% of A records and 26\% of B records. Second, the no-alarm fallback should remain an explicit, auditable state rather than being replaced by the largest finite threshold, since the latter can still alarm and has no certificate. Third, a target site that has local labels should recalibrate on its own patient blocks before adopting a source threshold, because the external rate and sensitivity can move in opposite directions; our oracle split shows that this recovers part of the sensitivity but not without a higher alarm load, and a label-free version remains an open deployment problem. Fourth, the threshold grid and the realised-rate report should be published together, since a grid floor can make high budgets look identical and hide the fact that the constraint is no longer binding. These practices are operational rather than architectural, and they are compatible with the alert-stewardship literature that treats interruptive alerts as a monitored resource \citep{khan2026qualitative,ray2026alert,jaidev2026sepsis}.

The study has several limitations that bound its claims. It is retrospective and uses a single public challenge cohort with two sites; the subject-disjoint split uses the official one-file-per-subject structure, but the release cannot link separate readmissions, so the independent-subject assumption remains an assumption. The label is the challenge's sepsis definition, not a adjudicated clinical deterioration event, and left-censored records are excluded rather than imputed. The score models are deliberately simple reference learners, LightGBM is not tuned, and the GRU is not a state-of-the-art architecture. The partial NEWS, SIRS and qSOFA baselines omit components that the 2019 release does not contain, so they are research comparators rather than certified clinical instruments. No label-free target-site adaptation, subgroup analysis, observation-sparsification study or clinical outcome evaluation was performed; the target-site recalibration and label-budget analyses are oracle diagnostics that use target labels and are not deployment procedures. ECE is low in absolute terms but is computed on a low-prevalence timestep problem and should not be read as clinical calibration. The primary external evaluation has no B calibration set by design, so it measures transport failure rather than the best achievable B performance. The three-seed E2 comparison is complete for all three backbones, but the target-site diagnostic for the GRU is single-seed, the label-budget curve is conditional on one target split and one policy family, and the study does not estimate GPU-hours or alert-fatigue outcomes.

These limitations do not overturn the methodological result, but they change what should be claimed. The paper does not show that budgeted calibration improves patient outcomes, reduces mortality or is preferred by clinicians. It shows that a stateful policy can be calibrated against a patient-level alarm budget with a finite-sample bound, that the resulting operating curve is measurable on a public cohort, and that the same threshold does not transport unchanged across hospitals. That is a narrower contribution than a clinical-effectiveness claim, but it is a more honest one, and it identifies the target-calibration experiment that a deployment study would need to run next.
